\documentclass[Total.tex]{subfiles}
\begin{document}
	\chapter{Holomorpihc Functions in Severial Variables}
		\section{Fundamentals}
			\subsection{Analysis and Topologies}
				Recall：
				\begin{itemize}
					\item Power Series $\CC[[z_1,\dots,z_n]]$ with elements $\sum a_{i_1,\dots,i_n}z_1^{i_1}\dots z_n^{i_n}$, $a_{i_1,\dots,i_n}\in \CC$;
					\item Convergence of Power Series;
					\item Uniformly Convergence to function $f:\CC^n\supseteq M\rightarrow \CC$,
					\begin{gather*}
						|\sum a_{i_1,\dots,i_n}z_1^{i_1}\dots z_n^{i_n}-f(z)|<\epsilon, \forall z\in M
					\end{gather*}
					\item Holomorphic Function $\Leftrightarrow$ Analytic
					
					According to Cauchy's theorem in multivariables, we obtain the equivalence.
					\begin{gather*}
						f\mbox{ is holomorphic}\Leftarrow f\mbox{ is analytic:} \forall z_0\in M,\exists U(z_0),
						\sum a_v(z-z_0)^v \mbox{converges on }U(z_0)\mbox{ to }f(z)
					\end{gather*}
					\item Absolute Space: A pointed set 
					\begin{gather*}
						V=\{(z_1,\dots,z_n)\in \CC^n:|z_i|<r_i\}
					\end{gather*}

					With the modular mapping on each component,\begin{gather*}
						\tau:\CC^n\rightarrow V\qquad (z_1,\dots,z_n)\mapsto (|z_1|,\dots,|z_n|)
					\end{gather*}
				\end{itemize}

				\subsubsection{Polycylinder}
				\begin{Def}
					Let	$\tau\in V$ be an element in absolute space,$z_0\in \CC^n$.
					
					then define the \textbf{polycylinder about $z_0$ with radius $\tau$} as
					\begin{gather*}
						\Delta(z_0,\tau)=\{z\in \CC^n:|z_i-z_{0,i}|<\tau_i, i=1,\dots,n\}
					\end{gather*}

					And, $\partial \Delta(z_0,\tau)$ is the \textbf{distinguished boundary} of $\Delta(z_0,\tau)$.
				\end{Def}

				It is easy to see that the geometric shape of polycylinder is actually a \textbf{product of discs} in $\CC^n$. Thus, the boundary is a product of circles, actually 'torus' in $\CC^n$.

				Now, take the following notion
				\begin{gather*}
					\odot{\CC^n}:='CC^n-\{0\}^n\qquad P_\tau=\Delta(0,\tau)
				\end{gather*}

			\subsection{Rheinhardt Domain}
				\subsubsection{Torus Action}
				\subsubsection{Definition of Rheinhardt Domain}
				\begin{Def}
					Let $D\subseteq \CC^n$ be a domain, then $D$ is called a \textbf{Rheinhardt Domain} if for any $z\in D$, the polycylinder $\Delta(0,|z|)\subseteq D$.
				\end{Def}
			\subsection{Complex Differential and Cauchy-Riemann Equation}
				\begin{Def}
					内容...
				\end{Def}
			\subsection{Some of Conclusions}
				\begin{lemma}
					内容...
				\end{lemma}
		\section{Theorems}
		\subsection{Cauchy's Formula}
			\begin{proposition}
				(Cauchy) If $f$ is a function defined in a neighborhood $U$ of a \underline{closed polydisc} $\bar{\Delta}(a,r)=\prod D(a_n,r_n)$ and holomorphic in(contains) a neighborhood $U$ and holomorphic in each variable $z_i$ at each point $U$, then
				\begin{gather*}
					f(z)=(\frac{1}{2\pi i})^n \int_{|\xi_n-a_n|=r_n}\dots \int_{|\xi_1-a_1|=r_1} \frac{f(\xi_1,\dots,\xi_n)}{\prod (\xi_i-z_i)} d\xi_i,z\in \bar{\Delta}(a,r)
				\end{gather*}
			\end{proposition}
			\begin{proof}
				Omitted,from Cauchy's formula in one variable.
			\end{proof}
			Note that:
			\begin{itemize}
				\item 
				\item 
			\end{itemize}
			
			\begin{proposition}
				(Osgood's Lemma)
			\end{proposition}
			\begin{proof}
				内容...
			\end{proof}
			
			\begin{proposition}
				(Cauchy's Inequality)
			\end{proposition}
		\subsection{Hartog's Theorem}
		\begin{proposition}
			(Jensen's Inequality)
			\end{proposition}
		\begin{proof}
			内容...
		\end{proof}
		\begin{proposition}
			内容...
		\end{proposition}
		\begin{proof}
			内容...
		\end{proof}
		
		\subsection{Convergence Theorem}
		\section{Domains of Holomorphy}
\end{document}